\scanpage{134}

% L1–2, title
\uline{Fibrés en drapeaux associés à}\deleted{ux}\inserted{ une} \uline{formes quadratique\unsure{s} régulière.}

% L3–5
$(E,q)$, avec $E$ de rg $N=2n+1$ (cas impair) ou $2n$, $n\geq1$ \inserted{donc $N\geq2$.} Soit $H\subset[1,n]\cap\mathbb Z$, \inserted{donc $H=\{i_1,\ldots,i_\nu\}$, où $1\leq i_{\unsure{1|j}}<\ldots<i_\nu\leq n$} \deleted{On pose} $\nu=\operatorname{card}H$

% L6–8
\[
 \begin{aligned}
 \Delta_H(E,q)&=\text{ens. des drapeaux de type }H\text{ de }E\\
 &E_1\subset E_2\ldots\subset E_\nu\quad(\nu=\operatorname{card} H,\ H=\{\operatorname{rg} E_1,\ldots,\operatorname{rg} E_\nu\}\\
 &\text{tels que }\deleted{\text{les}}\ E_\nu\text{ soit tot.}^{\mathrm t}\text{ isotrope.}
 \end{aligned}
\]

% L9
\[
 \underline\Delta_H(E,q):\mathrm{Sch}_{/S}\longrightarrow\mathrm{Ens}\qquad S'\longmapsto\Delta_H(E_{S'},q_{S'})
\]

% L10–12
\uline{Théorème 1} a) $\underline\Delta_H(E,q)\hookrightarrow\underline{\mathrm{Drap}}_H(E)$ est \deleted{représentable par} \inserted{une} immersion fermée -- en particulier $\underline\Delta_H(E,q)$ est représentable, et projectif sur $S$.

% L13–17
b) $\underline\Delta_H(E,q)$ est lisse sur $S$. Il est à fibres géométriquement $0$-connexes, sauf pour \deleted{$i_\nu=n$,} $\operatorname{rg}E$ pair $=2n$, $i_\nu=n$, auquel cas \deleted{\unsure{il}} ses fibres géom. ont exactement deux composantes connexes.

% L13–22, slanting left margin
% slanted
\margin{Donner aussi la structure de $\underline\Delta_H\to\underline\Delta_K$ pour $K\subset H$\ldots\ (Si $K=\varnothing$, on trouve $\underline\Delta_H\to S$\ldots) $\underline\Delta_H\to\underline\Delta_K$ est lisse, projectif, ses fibres géom sont $0$-conn. sauf si \deleted{\unsure{deux}} $\operatorname{rg} E$ est pair, $i_\nu=\operatorname{Max} H=n=\frac12\operatorname{rg} E$ \unsure{$\notin K$}, auquel cas les fibres géom. ont $2$ comp. conn.}

% L18–19
\uline{Dém} Le fait que $\underline\Delta_H(E,q)\to\mathrm{Drap}_H(E)$ une imm. fermée est immédiat, \unsure{..} .

% L20–22
\deleted{\unsure{Car si} $(E,q)$ est un}\inserted{Pour la lissité de} $\underline\Delta_H(E,q)$, on utilise le critère fonctoriel de lissité -- on est ramené à prouver que si $S_1\subset S$ est défini par $\mathcal J$ Idéal qu. coh. de carré nul,

% lower left
% archive sheet number: 133
